\documentclass{article}
\usepackage[UTF8]{ctex}
\usepackage{amsmath,amssymb,amsfonts}
\usepackage{enumitem}
\usepackage{geometry}
\geometry{a4paper, margin=1in}
% === 错题本排版优化（不参与题目标识） ===
\setlength{\parindent}{0pt}
\setlength{\parskip}{0.35em plus 0.10em minus 0.08em}
\setlength{\emergencystretch}{3em}
\setlist{itemsep=0.25em, topsep=0.40em, parsep=0pt}
\allowdisplaybreaks[2]
% === 错题本排版优化结束 ===
\begin{document}
\section*{厦门大学2026年数学分析真题}
\section*{115．厦门大学2026年数学分析试卷}

\subsection{题目1}
1．（15 分）设 \(a_{n}>0(n=1,2, \cdots)\) ，级数 \(\sum_{n=1}^{\infty} a_{n}\) 发散，证明：\(\sum_{n=1}^{\infty} \frac{a_{n}}{S_{n}}\) 发散，其中 \(S_{n}=a_{1}+a_{2}+\cdots+a_{n}\) 。

\subsection{题目2}
2．（15 分）设 \(x_{n}=1+\frac{1}{2}+\cdots+\frac{1}{n}-\ln (n+1)(n=1,2, \cdots)\) ，证明： \(\lim _{n \rightarrow \infty} x_{n}\) 存在．

\subsection{题目3}
3．（20 分）设 \(f(x)\) 是定义在 \([0,1]\) 上的连续可微函数，且满足 \(f(0)=0, f(1)=1\) ，证明：

\[
\int_{0}^{1}\left|f(x)-f^{\prime}(x)\right| \mathrm{d} x \geq \frac{1}{e}
\]

\subsection{题目4}
4．（20 分）设 \(f(x)\) 在 \((-\infty,+\infty)\) 上有三阶连续导数，且 \(f(x), f^{\prime \prime \prime}(x)\) 在 \((-\infty,+\infty)\) 上有界，证明： \(f^{\prime}(x), f^{\prime \prime}(x)\) 在 \((-\infty,+\infty)\) 上有界．

\subsection{题目5}
5．（20 分）设 \(f(x), g(x)\) 是定义在 \([0,+\infty)\) 上的非负连续函数，满足

\[
f(x) \leq \int_{0}^{x} f(s) g(s) \mathrm{d} s, \forall x \geq 0
\]

证明：\(f(x) \equiv 0, x \in[0,+\infty)\) ．

\subsection{题目6}
6．（20 分）记 \(\mathbb{R}^{3}\) 中的点 \(x=\left(x_{1}, x_{2}, x_{3}\right), \mathbb{R}^{3}\) 中以原点为圆心，\(r>0\) 为半径的开球记为 \(B_{r}\) ，其边界为\(\partial B_{r}\) ．设 \(f: \mathbb{R}^{3} \rightarrow \mathbb{R}\) 具有连续的偏导数，且 \(f(x)=0, \forall x \in \partial B_{r}\) ，证明：

\[
\lim _{\varepsilon \rightarrow 0^{+}} \iiint_{B_{r} \backslash B_{\varepsilon}} \frac{\sum_{i=1}^{3} x_{i} \frac{\partial f}{\partial x_{i}}(\mathbf{x})}{|\mathbf{x}|^{3}} \mathrm{~d} V=-4 \pi f(0)
\]

\subsection{题目7}
7．（20 分）设 \(V\) 是由球面 \(x^{2}+y^{2}+z^{2}=2 a z\) 与 \(x^{2}+y^{2}+z^{2}=a z\) 所围立体（ \(a>0\) ），计算 \(\iiint_{V} z \mathrm{~d} x \mathrm{~d} y \mathrm{~d} z\) ．

\subsection{题目8}
8．（20 分）设 \(f(x)\) 在 \([-\pi, \pi]\) 上连续，\(a_{n}(n=0,1, \cdots), b_{n}(n=1,2, \cdots)\) 为 \(f\) 的 Fourier 系数．证明：级数 \(\sum_{n=1}^{\infty} \frac{a_{n}}{n}\) 和 \(\sum_{n=1}^{\infty} \frac{b_{n}}{n}\) 均收敛。
\end{document}
